\section{Hardware Requirements}\label{sec:hardware}

\paragraph{Google datacenters.}
The SRE book, which carries a 2017 copyright, says that most of Google's compute resources are in Google-designed datacenters with proprietary power distribution, cooling, networking and compute hardware \cite{sre-book}. Google designs its server boards and networking equipment \cite{docs-cloud-google-docs-security-infrastructure-design} and custom builds servers exclusively for its data centers \cite{datacenters}. Tens of machines sit in a rack, and a datacenter building usually houses several clusters \cite{sre-book}. The SRE book's main text says the compute hardware is the same across the board, but a footnote says some datacenters end up with multiple generations \cite{sre-book}, and the load-balancing chapter agrees that machines in one datacenter are not necessarily the same \cite{sre-book}. Hardware failure is routine: in a single cluster in a typical year, thousands of machines fail and thousands of hard disks break \cite{sre-book}. The SRE book also says that Google's Jupiter network connects hundreds of Google-built switches in a Clos fabric and supports 1.3~Pbps of bisection bandwidth in its largest configuration \cite{sre-book}. The B4 backbone connects datacenters to each other \cite{sre-book}. The lowest layer of the storage stack uses both spinning disks and flash \cite{sre-book}. Besides its own datacenters, Google has racks of proxy and cache machines in colocation facilities \cite{workbook}.

\paragraph{Firmware and processors.}
Google manages the firmware of its machines and network cards with the same systems it uses for host software, and names Arm and x86 CPUs among its environments \cite{bsrs}. It designs custom chips, including the Titan hardware security chip, which serve as hardware roots of trust \cite{docs-cloud-google-docs-security-infrastructure-design}.

\paragraph{Sizing numbers are exercise assumptions.}
The SRE Workbook's design exercise assumes a standard machine of 16 cores, 64~GB of RAM and 1~Gbps of network throughput, and a common 4~TB disk sustaining about 200 operations per second \cite{workbook}. These are assumptions of the exercise, not a statement of Google's hardware.

\paragraph{Hardware for practice (recommendation).}
Debian's installation guide lists 512~MB of RAM as the minimum and 1~GB as recommended, with 4~GB of disk, for an install without a desktop \cite{debian-releases-stable-amd64-ch03s04-en-html}. The Ubuntu Server basic-installation tutorial recommends at least 2~GB of RAM and 5~GB of disk \cite{ubuntu-server-docs-tutorial-basic-installation}. Minikube needs 2~CPUs, 2~GB of free memory and 20~GB of free disk \cite{minikube-sigs-k8s-docs-start}. From these minimums we estimate, as an assumption and not a source's advice, that a personal laptop or desktop with about 8~GB of RAM can run either minikube or two or three small Debian virtual machines without a desktop, and that 16~GB runs both at once. Server-class hardware is not needed.
